% TOP_PROBLEM: {"id": 7600006, "title": "Shelah's categoricity conjecture for $L_{\\omega_1,\\omega}$", "category_id": 18, "category": {"id": 18, "name": "logic", "display_name": "Logic", "description": "Problems in mathematical logic, model theory, proof theory, and finite model theory.", "slug": "logic", "order_index": 18, "created_at": "2026-07-31T15:26:25.671Z"}, "status": "partially_solved"}
% FIRST_POSTED: 2026-10-06
% ATTEMPT_STATUS: unresolved
% OPENAI_CLAIM: {"schema_version":1,"source_repo":"https://github.com/openai/math","source_commit":"adc7f1241b42e322a6451854ab7e4b4c146bf78a","release_date":"2026-10-06","match":"partial","resolution":"partial","independently_reviewed":false,"second_pass":"n/a","families":[{"family":"240","title":"Shelah's eventual categoricity and the prescribed-threshold obstruction","subject":"Mathematical logic","manuscripts":[{"title":"Eventual categoricity for abstract elementary classes","pdf_path":"preprints/Eventual-Categoricity-for-Abstract-Elementary-Classes-September-24-2026/paper.pdf","readme_path":"preprints/Eventual-Categoricity-for-Abstract-Elementary-Classes-September-24-2026/README.md","date":"2026-09-24","theorem_ref":"main:uniform (sections/01-introduction.tex:45); main:qualitative (sections/01-introduction.tex:75)","statement_sources":[{"path":"preprints/Eventual-Categoricity-for-Abstract-Elementary-Classes-September-24-2026/build/sections/01-introduction.tex","line":45,"label":"main:uniform"},{"path":"preprints/Eventual-Categoricity-for-Abstract-Elementary-Classes-September-24-2026/build/sections/01-introduction.tex","line":75,"label":"main:qualitative"}]}],"lean":{"doc_path":"lean/docs/240.md","comparators":[],"declarations":[],"covers_main_theorem":false},"reasoning_trace":null}]}
\documentclass[11pt]{article}
\usepackage[margin=25mm]{geometry}
\usepackage{fontspec}
\setmainfont{Latin Modern Roman}
\usepackage{amsmath,amssymb,mathtools}
\usepackage{unicode-math}
\setmathfont{Latin Modern Math}
\usepackage{enumitem}
\usepackage{xurl,hyperref,bookmark}
\hypersetup{colorlinks=true,urlcolor=blue}
\setcounter{secnumdepth}{0}
\setlength{\parindent}{0pt}
\setlength{\parskip}{5pt}
\setlength{\emergencystretch}{3em}
\begin{document}
\section*{Shelah's categoricity conjecture for \$L\_\{\textbackslash{}omega\_1,\textbackslash{}omega\}\$}

\subsection{OpenAI's claim}
OpenAI presents the following claims relevant to this question, with manuscript numbers following paper order in Sources. In manuscript 1 (family 240; 2026-09-24), the asserted result is that in ZFC, each bound on L\"owenheim--Skolem number has a uniform threshold above which categoricity in one cardinal implies categoricity throughout the same tail.

\subsection{Mathematical statement of the claimed result}
\paragraph{Manuscript 1.}
For every infinite cardinal $\lambda$, the claim gives a uniform cardinal threshold $\mu(\lambda)\ge\lambda$ for all abstract elementary classes $K$ of L\"owenheim--Skolem number at most $\lambda$:
\[
 (\exists\kappa\ge\mu(\lambda):I(K,\kappa)=1)
 \Longrightarrow(\forall\kappa'\ge\mu(\lambda):I(K,\kappa')=1).
\]
Here $I(K,\kappa)$ is the number of model isomorphism classes of cardinality $\kappa$. The same threshold occurs in the premise and conclusion. A second formulation asserts that categoricity in unboundedly many cardinals forces categoricity on a tail, without adding amalgamation or tameness assumptions.

\subsection{Outline of the argument}
For manuscript 1 (family 240; 2026-09-24), fixed order-indexed presentations encode short model orbits. A direct reversal theorem for strong maps preserves bounded tuple records; elimination and stabilization of strong-map classes then create the stable diagrams and closed-level amalgamation needed for tail transfer.

\subsection{Scope relative to this problem}
The eventual-categoricity paper supplies a uniform sufficiently-large threshold but does not identify it with the particular Hanf or beth threshold in this statement. Transfer at cardinals between that specified bound and the larger threshold is consequently not supplied.

\subsection{Formal verification}
No Lean formalization is listed at this snapshot for manuscript 1 (family 240; 2026-09-24).

This site has not executed the Lean code. Described statement coverage is not an independent proof check.

\subsection{Status caveat}
The manuscripts were produced by an internal OpenAI model. OpenAI cautions that results without formalization can contain mistakes. This site has not independently checked the proofs; the label records the scope of OpenAI's claim.

\subsection{Sources}
\begin{itemize}
\item OpenAI, \emph{Eventual categoricity for abstract elementary classes}, 2026-09-24; citation key \texttt{OAI:Eventual-Categoricity-for-Abstract-Elementary-Classes-September-24-2026}. \href{https://github.com/openai/math/blob/adc7f1241b42e322a6451854ab7e4b4c146bf78a/preprints/Eventual-Categoricity-for-Abstract-Elementary-Classes-September-24-2026/paper.pdf}{PDF (GitHub)}; \href{https://github.com/openai/math/raw/adc7f1241b42e322a6451854ab7e4b4c146bf78a/preprints/Eventual-Categoricity-for-Abstract-Elementary-Classes-September-24-2026/paper.pdf}{PDF (direct)}; \href{https://github.com/openai/math/blob/adc7f1241b42e322a6451854ab7e4b4c146bf78a/preprints/Eventual-Categoricity-for-Abstract-Elementary-Classes-September-24-2026/README.md}{Manuscript page and citation}; \href{https://github.com/openai/math/blob/main/preprints/Eventual-Categoricity-for-Abstract-Elementary-Classes-September-24-2026/paper.pdf}{Latest version}; Statement source: \href{https://github.com/openai/math/blob/adc7f1241b42e322a6451854ab7e4b4c146bf78a/preprints/Eventual-Categoricity-for-Abstract-Elementary-Classes-September-24-2026/build/sections/01-introduction.tex\#L45}{main:uniform}; Statement source: \href{https://github.com/openai/math/blob/adc7f1241b42e322a6451854ab7e4b4c146bf78a/preprints/Eventual-Categoricity-for-Abstract-Elementary-Classes-September-24-2026/build/sections/01-introduction.tex\#L75}{main:qualitative}.
\item Family 240: \href{https://github.com/openai/math/blob/adc7f1241b42e322a6451854ab7e4b4c146bf78a/lean/docs/240.md}{Lean scope}.
\item \href{https://github.com/openai/math/blob/adc7f1241b42e322a6451854ab7e4b4c146bf78a/CONTENTS.md}{OpenAI catalogue entry}.
\end{itemize}
\end{document}
